% Section 7.7, from the summary table of lectures/L07/appendix/a_theory.md and from
% lectures/L07/README.md: the objectives, the evaluation questions and the next lesson.

\section{Sammanfattning}
\label{c7:sec:sammanfattning}

\begin{booktable}{@{}p{12.5em}L@{}}
  \thead{Begrepp} & \thead{Formel} \\
  \midrule
  Absolutbelopp & $\abs{\mathbf{u}} = \sqrt{u_x^2 + u_y^2}$ \\
  Vinkel mot $x$-axeln & $v = \arctan(u_y / u_x)$ + kvadrantkorrigering \\
  Skalärprodukt & $\mathbf{u} \cdot \mathbf{v} = u_x v_x + u_y v_y$ \\
  Vinkel mellan vektorer &
    $\cos \alpha = \dfrac{\mathbf{u} \cdot \mathbf{v}}{\abs{\mathbf{u}}\abs{\mathbf{v}}}$ \\
  Addition & $\mathbf{u} + \mathbf{v} = (u_x + v_x;\, u_y + v_y)$ \\
  Skalärmultiplikation & $k\mathbf{u} = (ku_x;\, ku_y)$ \\
\end{booktable}

\needspace{6\baselineskip}
\subsection*{Det här ska du kunna}
Efter det här kapitlet ska du:
\begin{itemize}
  \item Förstå begreppet vektor och hur en vektor representeras matematiskt.
  \item Kunna genomföra enkla beräkningar med vektorer.
\end{itemize}

\testadigsjalv
\begin{enumerate}
  \item En vektor är $\mathbf{u} = (3, -4)$. Beräkna absolutbeloppet $\abs{\mathbf{u}}$ och
    vinkeln mot $x$-axeln.
  \item Beräkna $\mathbf{w} = 2\mathbf{u} - \mathbf{v}$ för $\mathbf{u} = (2, 1)$ och
    $\mathbf{v} = (-1, 3)$.
\end{enumerate}

\needspace{6\baselineskip}
\subsection*{Nästa kapitel}
\Kapref{8} handlar om funktioner.
